\subsection{正交群中的对称。}

除非另有明确说明，在本小节中我们假设 A 是特征 ≠ 2 的交换域，且 Φ 是与 E 上一个非退化二次型 Q 相伴的对称双线性型。回忆一下，对 x ∈ E 有 Φ(x, x) = 2Q(x) （§ 3，第 4 小节）。

设 H 是 E 中一个非迷向超平面，u 是关于 H 的对称（第 3 小节）。设 \(a \neq 0\) 是与 H 正交的一个向量；依假设有 \(u(a) = -a\) 。每个向量 \(x \in E\) 都可以以唯一的方式写成 \(x = \lambda a + y\) 的形式，其中 \(\lambda \in A\) 且 \(y \in H\) ；由于 \(a\) 与 \(y\) 正交，有 \(\Phi(x, a) = \lambda \Phi(a, a)\) ，于是，由于 \(a\) 非迷向（ \(\S 4\) ，第 1 小节，命题 1 的推论），\(\lambda = \Phi(x, a) \Phi(a, a)^{-1}\) 。在这些条件下，有

\[
u (x) = \lambda u (a) + u (y) = - \lambda a + y = x - 2 \lambda a,
\]

由此得

\[
u (x) = x - 2 \Phi (x, a) \Phi (a, a) ^ {- 1}. a = x - \Phi (x, a) Q (a) ^ {- 1}. a.\tag{6}
\]

注意，当 A 是特征 2 的体且 a 是 E 的一个非奇异向量时，(6) 的最后一项仍保留其意义；我们立即验证，此时仍有 Q(u(x)) = Q(x) （对一切 x ∈ E），换句话说，u ∈ O(Q)。我们也称如此定义的对合 u 是关于与 a 正交的超平面的对称（参见习题 28）。

命题 5. --- 假设向量空间 E 是有限维 n 的。此时与 Q 相伴的正交群 \(\mathbf{O}(\mathrm{Q})\) 由关于 E 的非迷向超平面的诸对称生成。

\(n=0\) 时命题是显然的，我们对 \(n\) 作归纳推理。设 \(u\) 是 E 的一个正交变换，\(x\) 是 E 的一个非迷向向量（引理 1）；我们区分三种情形：

\begin{enumerate}
\def\labelenumi{\alph{enumi})}
\item
  先假设 \(u(x) = x\) 。则与 \(x\) 正交的超平面 H 是非迷向的，并且我们有 \(u(\mathrm{H}) = \mathrm{H}\) 。因此限制 \(u'\) （\(u\) 在 H 上的）属于正交群 \(\mathbf{O}(Q')\) （与 Q 在 H 上的限制 \(Q'\) 相伴的）。由于 \(Q'\) 非退化，归纳假设蕴含 \(u' = v_1' \ldots v_m'\) ，其中 \(v_i'\) 是关于 H 的一个超平面 \(L_i\) 的对称。于是自同态 \(v_i\) （E 的、扩张 \(v_i'\) 且使 \(v_i(x) = x\) 的）是关于 E 的超平面 \(Ax + L_i\) 的对称。我们显然有 \(u = v_1v_2 \ldots v_m\) 。
\item
  其次假设 \(u(x) = -x\) 。若以 \(s\) 表示关于与 \(x\) 正交的超平面 H 的对称，并令 \(\nu = su\) ，则有 \(\nu(x) = x\) ，于是化归到情形 \(a\) )。
\item
  最后转到一般情形，令 \(y = u(x)\) ，于是 \(Q(y) = Q(x)\) 。在这些条件下，向量 \(x - y\) 与 \(x + y\) 不可能都是迷向的，因为从关系 \(Q(x - y) = 0\) 与 \(Q(x + y) = 0\) ，逐项相加就将得出 \(2(Q(x) + Q(y)) = 0\) （§ 3，第 4 小节，定义 2），从而 \(4Q(x) = 0\) ，与假设相反。例如设 \(a = x - y\) 不是迷向的；此时我们有
\end{enumerate}

\[
\Phi (y, a) = \mathrm{Q} (y + a) - \mathrm{Q} (y) - \mathrm{Q} (a) = \mathrm{Q} (x) - \mathrm{Q} (y) - \mathrm{Q} (a) = - \mathrm{Q} (a);
\]

因此，若以 \(s\) 表示关于与 \(a\) 正交的超平面的对称，则公式 (6) 证明 \(s(y) = y + a = x\) ；令 \(\nu = su\) ，就有 \(\nu(x) = x\) ，于是化归到情形 \(a\) 。若 \(a = x - y\) 迷向而 \(b = x + y\) 不迷向；同样可以看出化归到情形 \(b\) 。

